The full benchmark separates the benefit of compressing a dense objective from
the effect of changing that objective. At tolerance 0.01, plain low rank has
lower global overlap than the quality-preset graph comparator on
\BenchPlainGraphGlobalWins{} of \BenchPlainGraphGlobalN{} molecular datasets
(Figure~\ref{fig:benchmark-global}). The combined setting reduces local leakage
relative to plain low rank on \BenchCombinedPlainLocalWins{} of
\BenchCombinedPlainLocalN{} datasets, while raising the ordinary dense cut.
Its local score is lower than the graph comparator on
\BenchCombinedGraphLocalWins{} of \BenchCombinedGraphLocalN{} datasets
(Figure~\ref{fig:benchmark-local}). The mean combined-minus-graph local
change is \BenchCombinedGraphLocalChange{}, reflecting the differing
dataset effects. These comparisons expose the choice the
framework offers: the plain factor directly targets distributed overlap;
power and grouping spend some of that global advantage on local separation.
The accompanying machine-readable tables include both tolerances, all reference methods, realized
counts and group relaxations. Paired dataset effects, bootstrap intervals and
Holm-adjusted sign tests are supplied in the supplement.

\begin{figure}[!htbp]
\centering\includegraphics[width=6.5in]{revision/benchmark_global.pdf}
\caption{Global overlap on eleven full molecular datasets at size tolerance
0.01, ordered by dataset size. Points are means of three split seeds; bars show
their standard deviations. The dotted line is the count-matched random
expectation. All methods retain the same input entities. The graph comparator
uses exact top-32 construction, an ordinary pairwise Graph API, the quality
preset and CUT objective, with lower/upper count repair. The supplement reports
all methods, actual counts, degree sensitivity and the two materials checks.}
\label{fig:benchmark-global}
\end{figure}

\begin{figure}[!htbp]
\centering\includegraphics[width=6.5in]{revision/benchmark_local.pdf}
\caption{Test-only nearest-training-neighbour leakage for the same partitions
and scoring protocol as Figure~\ref{fig:benchmark-global}. Ratios are normalized
by count-preserving random permutations. Bars show split-seed standard
deviations, truncated at zero where necessary. Controlling total overlap and
close-neighbour contamination gives different method orderings.}
\label{fig:benchmark-local}
\end{figure}

DataSAIL provides a closer comparison for the dense objective. Its saved
assignments cover seven datasets, with source outcomes classified separately
from missing attempts and resource-policy exclusions. At their exact native
counts, plain low rank has lower global overlap on
\BenchPlainDataSailGlobalWins{} of \BenchPlainDataSailGlobalN{} datasets.
The combined setting has lower global overlap on
\BenchCombinedDataSailGlobalWins{} of \BenchCombinedDataSailGlobalN{} and lower
local leakage on \BenchCombinedDataSailLocalWins{}, with
\BenchCombinedDataSailLocalTies{} tie. Thus compression provides a practical
dense-cut objective, while the local setting remains a choice among competing
criteria. Reused assignments are rescored with the common protocol and supply
no new runtime.

The published local and cluster methods clarify the retention contract.
Lo-Hi satisfies its 0.4 cross-fold threshold on all seven eligible datasets,
retaining between 98.0\% and 99.9\% of their molecules. Low-rank controls match
each retained subset and its exact counts; their capped groups offer soft local
control rather than that hard threshold guarantee. In the transferred UMAP
recipe, seven native cluster holdouts and their count-matched controls retain
every molecule. The supplement reports both studies separately. The FreeSolv native UMAP
failure and MUV memory-policy exclusion remain unscored.

\IfFileExists{runtime_numbers.tex}{
\subsection{Dense splitting has a distinct computational advantage}

The same-host comparison measures the public split call at exactly the native
DataSAIL counts. Across four cases, median combined-setting times range from
\RuntimeCombinedMinSeconds{} to \RuntimeCombinedMaxSeconds{} seconds;
DataSAIL takes \RuntimeDataSailMinSeconds{} to
\RuntimeDataSailMaxSeconds{} seconds, including distance construction,
clustering and its integer program (Figure~\ref{fig:runtime-matched}).
Local quality remains dataset-dependent: on BACE the combined mean local ratio
is \RuntimeBaceCombinedLocal{}, compared with
\RuntimeBaceDatasailLocal{} for DataSAIL. Speed and the desired separation
criterion must therefore be considered together. Solver outcomes, initialization
policy, individual timings and all exact counts appear in the supplement.

\begin{figure}[!htbp]
\centering\includegraphics[width=6.5in]{revision/runtime_matched.pdf}
\caption{CPU cost and local overlap on four same-host cases at exactly matched
DataSAIL fold counts. Colour denotes method and shape denotes dataset.
Low-rank points average three split seeds for local leakage and use median
split-call time; DataSAIL has one solve per case. The clock excludes feature
loading and scoring. DataSAIL's distance construction, clustering and solver
are included; low-rank calls exclude a 100-entity warm-up. Solver-reported
optima and time-limited feasible outcomes are distinguished in the supplement.}
\label{fig:runtime-matched}
\end{figure}
}{}
